\subsection{A second section with three residual families}
\label{sec:seven-three}

For the same forcing word and destination cycle take

\begin{equation}\label{eq:seven-new-section}
\begin{gathered}
 \mathcal S_* = \{(z,i):z=3m,\ m\ge1,\ i\in\{1,2,3,5\}\},
 \\
 W(m,i)=w_i m,\quad (w_1,w_2,w_3,w_5)=(4,7,4,5).
\end{gathered}
\end{equation}
The physical launch is still $s(m,i)=(3m-1+h_i,i)$, now with the
four indicated phases.  Write $P_*$ for the first strictly later visit
to $\mathcal S_*$.

\begin{theorem}[Three-family conditional capture]\label{thm:seven-three}
Every positive physical integer at every initial phase is captured by
\eqref{eq:seven-cycle} if and only if all three following statements hold:
\begin{align*}
 \mathrm{B1}:&\quad \operatorname{Cap}(s(2t,1)) &&(t\ge1),\\
 \mathrm{A3}:&\quad \operatorname{Cap}(s(4t+3,3)) &&(t\ge0),\\
 \mathrm{N03}:&\quad \operatorname{Cap}(s(8t+2,5)) &&(t\ge0).
\end{align*}
Every safe orbit enters $\mathcal S_*$ within eleven updates, and all
first returns from phases $1,2,3,5$ occur within $2,7,6,7$ updates,
respectively.  Outside the three families and five explicitly captured
states, $P_*$ strictly decreases the single natural-valued rank $W$.
The theorem does not prove any of the three universal capture hypotheses.
\end{theorem}

\paragraph{Uniform entry and bounded return.}
For entry, partition the phase-zero positive coordinate as
$z=1536q+r$, $0\le r<1536$, with $q\ge1$ for $r=0$ and $q\ge0$
otherwise.  The coefficient $1536=3\cdot2^9$ has more than enough
powers of two for five exact updates and retains a factor of three
throughout.  Consequently both parity and section membership are
constant on every propagated affine domain.  Independent physical
propagation gives the complete entry-time distribution
\[
 \begin{array}{c|rrrr}
 \text{first entry time}&1&2&3&5\\\hline
 \text{number of residue domains}&1024&384&64&64.
 \end{array}
\]
The counts sum to 1536.  Alignment from an arbitrary launch phase to
phase zero costs at most six updates, proving the stated bound of eleven.

For first returns, start with the entire domain $z=3m$ at each marked
phase.  Whenever the next parity is not fixed, split the free parameter
into its even and odd cases.  A growing branch $G$ and a shrinking branch
$S$ in the positive coordinate have the formulas
\[
 G_i(z)=\frac{3z+h_i}{2}\quad(z\equiv h_i\pmod2),\qquad
 S_i(z)=\frac{z+1-h_i}{2}\quad(z\not\equiv h_i\pmod2).
\]
These formulas describe branches of the prescribed $T/U$ schedule;
the branch is determined by the input parity. After $g$ growing updates,
the free coefficient is $3^{g+1}$: each parity split supplies the factor
of two consumed by that update. It is odd and divisible by three,
justifying the next split and the constant test for section membership.
Stop a domain at its
first later marked phase with $3\mid z$.  This procedure terminates
within the stated caps and produces the 43 domains in
Table~\ref{tab:seven-three-cylinders}.  At each source phase the
domains are pairwise disjoint and their dyadic widths sum to one.
Since all moduli are powers of two and bounded by the corresponding
$2^{L_i}$, this is an exhaustive finite partition of the positive
integers, not a density-only coverage claim.

As a separate reconstruction, partition $m$ modulo $2^{L_i}$ with
$(L_1,L_2,L_3,L_5)=(2,7,6,7)$, and propagate the original physical
recurrence, checking safety and the absence of an earlier section hit.
There are $4+128+64+128=324$ full-modulus input domains.  Each reproduces
the endpoint and time of its unique adaptive domain.  The two procedures
share the forcing and the section definition, but use different
coordinates and different partitions.

\begingroup
\small
\setlength{\tabcolsep}{7pt}
\begin{longtable}{rllllr}
\caption{The complete first-return partition for $\mathcal S_*$. All rows have
$t\ge0$, except $t\ge1$ when the displayed input is a multiple of $t$
with zero constant. The trace records every update, so its length is the
return time. The final column is $W_{\rm out}-W_{\rm in}$.}
\label{tab:seven-three-cylinders}\\
\toprule
$i$&input $m$&output $(m',j)$&trace&rank change&updates\\\midrule
\endfirsthead
\toprule
$i$&input $m$&output $(m',j)$&trace&rank change&updates\\\midrule
\endhead
\bottomrule
\endfoot
1&$2t$&$(3t,2)$&\texttt{G}&$13t$&1\\
1&$4t+1$&$(3t+1,3)$&\texttt{SG}&$-4t$&2\\
1&$4t+3$&$(t+1,3)$&\texttt{SS}&$-12t-8$&2\\
2&$2t$&$(3t,3)$&\texttt{G}&$-2t$&1\\
2&$8t+3$&$(9t+4,5)$&\texttt{SGG}&$-11t-1$&3\\
2&$8t+5$&$(3t+2,5)$&\texttt{SSG}&$-41t-25$&3\\
2&$8t+7$&$(3t+3,5)$&\texttt{SGS}&$-41t-34$&3\\
2&$64t+1$&$(9t+1,1)$&\texttt{SSSSGG}&$-412t-3$&6\\
2&$64t+17$&$(3t+1,1)$&\texttt{SSSSSG}&$-436t-115$&6\\
2&$64t+25$&$(9t+4,1)$&\texttt{SSSGGS}&$-412t-159$&6\\
2&$64t+33$&$(3t+2,1)$&\texttt{SSSSGS}&$-436t-223$&6\\
2&$64t+41$&$(9t+6,1)$&\texttt{SSSGSG}&$-412t-263$&6\\
2&$64t+49$&$(t+1,1)$&\texttt{SSSSSS}&$-444t-339$&6\\
2&$64t+57$&$(27t+25,1)$&\texttt{SSSGGG}&$-340t-299$&6\\
2&$128t+9$&$(9t+1,2)$&\texttt{SSSGSSG}&$-833t-56$&7\\
2&$128t+73$&$(3t+2,2)$&\texttt{SSSGSSS}&$-875t-497$&7\\
3&$4t$&$(3t,5)$&\texttt{SG}&$-t$&2\\
3&$4t+1$&$(3t+1,5)$&\texttt{GS}&$-t+1$&2\\
3&$4t+3$&$(9t+7,5)$&\texttt{GG}&$29t+23$&2\\
3&$32t+2$&$(9t+1,1)$&\texttt{SSGGS}&$-92t-4$&5\\
3&$32t+6$&$(3t+1,1)$&\texttt{SSSGS}&$-116t-20$&5\\
3&$32t+10$&$(9t+3,1)$&\texttt{SSGSG}&$-92t-28$&5\\
3&$32t+18$&$(27t+16,1)$&\texttt{SSGGG}&$-20t-8$&5\\
3&$32t+22$&$(9t+7,1)$&\texttt{SSSGG}&$-92t-60$&5\\
3&$32t+30$&$(3t+3,1)$&\texttt{SSSSG}&$-116t-108$&5\\
3&$64t+14$&$(3t+1,2)$&\texttt{SSSSSG}&$-235t-49$&6\\
3&$64t+26$&$(9t+4,2)$&\texttt{SSGSSG}&$-193t-76$&6\\
3&$64t+46$&$(t+1,2)$&\texttt{SSSSSS}&$-249t-177$&6\\
3&$64t+58$&$(3t+3,2)$&\texttt{SSGSSS}&$-235t-211$&6\\
5&$8t$&$(9t,1)$&\texttt{GSG}&$-4t$&3\\
5&$8t+2$&$(27t+7,1)$&\texttt{GGG}&$68t+18$&3\\
5&$8t+3$&$(9t+4,1)$&\texttt{SGG}&$-4t+1$&3\\
5&$8t+5$&$(3t+2,1)$&\texttt{SSG}&$-28t-17$&3\\
5&$8t+6$&$(9t+7,1)$&\texttt{GGS}&$-4t-2$&3\\
5&$8t+7$&$(3t+3,1)$&\texttt{SGS}&$-28t-23$&3\\
5&$16t+4$&$(3t+1,2)$&\texttt{GSSS}&$-59t-13$&4\\
5&$16t+9$&$(3t+2,2)$&\texttt{SSSG}&$-59t-31$&4\\
5&$16t+12$&$(9t+7,2)$&\texttt{GSSG}&$-17t-11$&4\\
5&$32t+17$&$(3t+2,3)$&\texttt{SSSSG}&$-148t-77$&5\\
5&$128t+1$&$(9t+1,5)$&\texttt{SSSSSGG}&$-595t$&7\\
5&$128t+33$&$(3t+1,5)$&\texttt{SSSSSSG}&$-625t-160$&7\\
5&$128t+65$&$(3t+2,5)$&\texttt{SSSSSGS}&$-625t-315$&7\\
5&$128t+97$&$(t+1,5)$&\texttt{SSSSSSS}&$-635t-480$&7\\
\end{longtable}
\endgroup


\paragraph{The finite exceptions and the induction.}
On a row $(Mt+r,i)\mapsto(At+b,j)$, the rank change is
\[
 W_{\rm out}-W_{\rm in}=(w_jA-w_iM)t+w_jb-w_ir.
\]
Exactly three rows have positive leading coefficient.  They are
\begin{center}
\begin{tabular}{llll}\toprule
Family&launch&first return&rank change\\\midrule
B1&$(2t,1)$&$(3t,2)$&$13t$\\
A3&$(4t+3,3)$&$(9t+7,5)$&$29t+23$\\
N03&$(8t+2,5)$&$(27t+7,1)$&$68t+18$\\\bottomrule
\end{tabular}
\end{center}
They grow on their whole legal domains.  Every other row has a negative
leading coefficient.  Solving its affine inequality on the entire legal
parameter ray leaves just
\[
 E=\{(1,1),(1,3),(1,5),(3,5),(5,3)\}
\]
as possible failures of strict descent.  Their capture times are
$4,2,0,14,16$ updates, respectively.  The first three have the physical
chains $(2,1)\to(1,2)\to(3,3)\to(4,4)\to(2,5)$ and its indicated
suffixes; $(2,5)$ is already on the destination cycle.  The longer two
chains, with phases advancing by one at each arrow, are
\begin{align*}
 (8,5):\quad&8,5,7,11,17,27,40,20,11,16,8,4,3,4,2,\\
 (15,3):\quad&15,22,11,18,8,4,2,2,0,0,1,1,2,1,3,4,2.
\end{align*}
Both end at phase five with physical value two.  The second chain
also illustrates why the convention at zero matters: it passes through
zero and continues, exactly as the forcing requires.

\begin{proof}[Proof of Theorem~\ref{thm:seven-three}]
The three families consist of positive physical states, so universal
capture implies their capture.  Conversely assume the three universal
statements.  If a safe orbit failed capture, uniform entry would give
a failed state in $\mathcal S_*$.  Choose such a state of least positive
integer rank $W$.  It is not in $E$, whose exact capture chains were
given above, and it is not in a residual family by hypothesis.
Its first return therefore strictly decreases $W$.  The return state
also fails capture, since its future is a suffix of the same orbit.
This contradicts minimality.  Every safe state is consequently captured.
\end{proof}

\paragraph{What changed relative to the four-family reduction.}
A3 is the old family S00.  At phase five, N03 is the subfamily
$m\equiv2\pmod8$ of the old even family S01.  The new section also
introduces B1 at phase one.  For the old S02 and S03 launches, it records
earlier endpoints
\[
 (32t+3,5)\xrightarrow{3}(36t+4,1),\qquad
 (128t+115,5)\xrightarrow{3}(144t+130,1).
\]
Their $W$-drops are $16t-1$ and $64t+55$.  The first is positive for
$t\ge1$; its remaining input $(3,5)$ belongs to $E$.  This explains
concretely how changing the section helps.  It does not prove capture
of those old families without the new hypotheses: a rank decrease
reduces an induction obligation but needs the common well-founded
argument to finish it.  Combining individual drops from $\rho$ and
$W$ without constructing one compatible order would not be valid.

\begin{proposition}[Three growing domains are optimal for this section]
\label{prop:seven-three-barrier}
Among positive phase-affine ranks
$R(m,i)=a_i m+b_i$ on $\mathcal S_*$, at least three of nine specified
disjoint first-return domains have a positive leading rank change.
The rank $W$ attains three growing domains in the complete 43-domain
partition.  The optimality assertion concerns this section and these
direct first-return comparisons.
\end{proposition}
\begin{proof}
Use the following three cycles of constraints, taking each labeled
edge from Table~\ref{tab:seven-three-cylinders}:
\[
 \begin{array}{c|c|c}
 \text{phase cycle}&\text{successive branch traces}&
       \text{product of slopes}\\\hline
 1\to2\to5\to1&G,\ SGG,\ GGS&243/128\\
 1\to3\to5\to1&SG,\ GS,\ GGG&243/128\\
 2\to3\to5\to2&G,\ GG,\ GSSG&243/128
 \end{array}
\]
Each edge is a different row of the disjoint partition.  If all three
leading rank changes in one cycle were nonpositive, multiplying
$a_jA\le a_iM$ around it would imply that the product of its slopes
$A/M$ is at most one.  Each displayed product is greater than one.
Thus each cycle contributes a different row with positive leading
coefficient; fixed offsets $b_i$ cannot cancel its growth for arbitrarily
large parameters.  The exact rank-change column proves attainment by
$W$.  These are cycles of inequalities on weights, not asserted
periodic trajectories satisfying every successive guard.
\end{proof}

\paragraph{Reproduction.}
\path{certificates/seven_step/three_family.py} regenerates
\path{three_family.json}, including the 43 affine rows, five finite
chains, independent physical reconstruction, entry times, and nine
obstruction witnesses. All 43 rows in \path{SevenRows.lean} agree with
the printed table after the parameter shift on zero-residue rows.

The Lean proof checks entry modulo 192 in at most six updates for
positive physical seeds at every phase. The eleven-update proof above
also covers safe zero states. Section~\ref{sec:formal-scope} gives the
formal three-family equivalence. The optimality proof above is written.
